\section{Small signed quartics and their arithmetic boundary}
\label{sec:signed}
Every variable in this section ranges over $\Z$. The new finite coordinate
and row transformations are applied to an authenticated signed arithmetic
source, not to a natural-number equation with its signs silently discarded.
This section first supplies the unbounded argument needed by that source;
then it proves the exact algebraic transformations used in the new endpoints.
The source files and each frozen transformation are part of the release.

\subsection{The unchanged signed arithmetic interface}
Write $t,F,B,V,y,x,N,q$ for the semantic coordinates before integral
centering. The unchanged signed Pell frame and semantic anchor identify
both coordinates of the principal Pell pairs before any low-digit sign
is used. This is the orientation-first interface retained from Round~17.
Its signed power lemma says that a framed power request at an integer base
$n\ge2$ forces its requested exponent $k$ to be nonnegative and its output
to equal $n^k$. Briefly, the classified signed Pell index gives the power
congruence; if $k<0$, multiplication by $n^{-k}$ and the strict frame bound
would force $mn^{-k}=1$ for an integer output $m$, which is impossible.
For base one and a known nonnegative index, the same congruence and frame
bound give output one. This argument does not assume the output is positive.

The retained ceiling signs and endpoint rows give
\begin{align}
 t&\ge16,&V&=(F^2+1)^t,\nonumber\\
 FB&=V+s_0,& B&=r_0+s_0+1,\qquad r_0,s_0\ge0,\nonumber\\
 y&=V(V-1),&x&=F+1+(y+1)u_0,\nonumber\\
 x&=1+[1+(t+1)y]u_t,&N&=B^x,\nonumber\\
 H_0&=1+\sum_{i=1}^{9}a_i,&q&=(BH_0)^x.
 \label{eq:signedfront}
\end{align}
The powers here describe the proved meanings of existing quadratic
certificates; exponentiation is not an operation in the final polynomial.
The $a_i$ are the actual encoded representatives. The derived constructor
label $h$ satisfies the retained ledger bound
\begin{equation}\label{eq:labelbound}
 0<h\le(n_*+1)^\ell\le(BH_0)^\ell,
 \qquad n_* =\max(1,a,b,c,d),\quad \ell=2^{22},
\end{equation}
once $a_i\ge0$ and $B\ge2$. The last inequality follows from
$n_*+1\le2H_0\le BH_0$, including the case $H_0=1$.

\subsection{An acyclic odd square bootstrap}
The odd-square source replaces the old ternary guard on $q-B-1$ by
\begin{equation}\label{eq:odd-square}
 q=(2z+1)^2.
\end{equation}
Only $q\ge1$ follows initially. The unchanged orientation argument uses
the square frame, not this missing sign. The request for $V$ has known
index $t\ge16$ and base at least one. Thus $V=(F^2+1)^t\ge1$.
The ceiling equations force $F\ge1$ and $B>F$, since
$FB\ge V\ge F^2+1$ while $B\ge1$. In particular $B\ge2$ and
$y>F+2$.

The genuine request $N=B^x$ now forces $x\ge0$. If $u_0\le-1$, the
first endpoint would give $x\le F-y<0$; hence $u_0\ge0$ and
$x\ge F+1\ge2$. The last endpoint then forces $u_t\ge1$, so
\begin{equation}\label{eq:endpointbound}
 x\ge2+(t+1)y.
\end{equation}
The numerator requests at base $q+1\ge2$ force all nine $a_i\ge0$;
the derived-label request and nonzero anchor give $h>0$. Consequently
$H_0\ge1$, and the existing request for $q$ yields
$q=(BH_0)^x\ge B^2>B$. The deleted output slack is therefore
$q-B-1\ge0$. Every original front sign has now been recovered, without
using a low-digit remainder sign or the conclusion $F=t!$.

\subsection{A uniform coefficient bound before factorial decoding}
The useful margin can be proved entirely before the factorial fixed point.
Since $V\ge2^t$, for $t\ge16$ we have
\[
 y\ge2^t(2^t-1)>\ell,\qquad y\ge t+2.
\]
Combining these integer inequalities with \eqref{eq:endpointbound} gives
\begin{equation}\label{eq:gaps}
 x\ge\ell t+4,\qquad x\ge t(t+1)+4.
\end{equation}
All requested digit positions have $k\le x$. For every coefficient below
the requested position, the following three estimates apply.

For an ordinary upper index $0\le a<H_0$ and $0\le j<k\le x$,
\[
 \binom aj\le\max(1,a)^j\le H_0^x\le q/16,
\]
because $B\ge2$ and $x\ge4$. For the derived label, $k=t+1$, so
\[
 \binom hj\le h^j\le(BH_0)^{\ell t}
 \le(BH_0)^{x-4}\le q/16\qquad(0\le j\le t).
\]
For the factorial upper index, $F^2+t\le B^{t+1}$: indeed
$F^2<B^2\le B^t$, $t\le2^t\le B^t$, and $2B^t\le B^{t+1}$.
Therefore
\[
 \binom{F^2+t}{j}\le(F^2+t)^t\le B^{t(t+1)}
 \le B^{x-4}\le q/16\qquad(0\le j\le t).
\]
These inequalities assume no factorial identity. Put
\[
 \rho(a,k)=\sum_{j=0}^{k-1}\binom ajq^j,\qquad D=q^k.
\]
For every needed pair $(a,k)$, the estimates give
\begin{equation}\label{eq:smalltail}
 0\le\rho(a,k)\le\frac{q(D-1)}{16(q-1)}
 \le\frac{D-1}{8}<\frac D8.
\end{equation}

\subsection{Recovering all remainder signs from shared budgets}
There are eleven remainders. The two with denominator $q^t$ are the
factorial digit remainder and encoded-index remainder; five have denominator
$q^{t+1}$, including the derived-label remainder; two stack remainders have
denominator $q^{y+1}$; and two strict quotient remainders have denominator
$q^x$. The old interval guards have the form
\[
 D=r+s+1,\qquad rs=u^2+v^2+w(w+1).
\]
The second expression is nonnegative at every integer tuple and represents
every nonnegative integer. To prove the latter assertion, apply the
three-square theorem to $4n+1$. Modulo four exactly one of its three square
summands is odd. Write it as $(2w+1)^2$ and the even squares as
$(2u)^2,(2v)^2$, then divide by four.

For the symmetric version, retain the linear interval rows and replace a
group's private product guards by
\begin{equation}\label{eq:symmetricbudget}
 D^2-4\sum_{i\in G}r_i^2-1=u^2+v^2+w(w+1).
\end{equation}
Every new zero has $|r_i|<D/2$. The unchanged numerator row gives
$r_i\equiv(q+1)^{a_i}\equiv\rho(a_i,k)\pmod D$. For the factorial
row its other part is $D(B+qH)$; for a range row it is $DNQ$. Both are
multiples of $D$ regardless of the signs of $H,Q$. Thus
\[
 |r_i-\rho(a_i,k)|<D/2+D/8<D,
\]
so the divisible difference is zero. All remainders are canonical and
nonnegative, and $s_i=D-1-r_i\ge0$ restores every old interval guard.
For the factorial digit, subtract its recovered tail and divide by $q^t$.
The quotient identity gives
\[
 B\equiv\binom{F^2+t}{t}\pmod q.
\]
The displayed coefficient bound also holds at $j=t$, so the binomial is
in $[0,q)$; the recovered front guard gives $0<B<q$. Thus
$B=\binom{F^2+t}{t}$ exactly. Only after this recovery is the inherited
factorial fixed-point theorem used to deduce $F=t!$. The unchanged CRT range and local Horn arguments
then produce an actual finite C222 derivation and its inherited ZFC
refutation interpretation.

Conversely, construct the inherited canonical witnesses from a finite
refutation. A group of at most five remainders obeys
$4\sum\rho_i^2<5D^2/16<D^2$. The nonnegative integer
$D^2-4\sum\rho_i^2-1$ therefore has a ternary representation, supplying
all new auxiliaries.

The adjacent levels merge by adding two product coordinates
$z_0=qr_0$, $z_1=qr_1$ and using one group at
$D'=q^{t+1}=qD$. Its seven entries are $z_0,z_1$ and the five higher-level
remainders. Soundness gives $|r_i|<D/2$ by division by $q>0$; completeness
uses $4\sum\rho_i^2<7D'^2/16<D'^2$. The two product rows are quadratic.
Replacing two ternary triples by one while adding two products saves one
coordinate. The alternative spherical version uses
$D^2-2\sum r_i^2=u^2+v^2+w(w+1)$; it gives
$|r_i|\le D/\sqrt2<3D/4$, still enough because $3D/4+D/8<D$.
Its seven-entry completeness bound is $2\sum\rho_i^2<7D^2/32<D^2$.
These are two different literal budget formulas, each with a complete
signed uniqueness argument.

\subsection{Completeness of the odd square restriction}
To construct \eqref{eq:odd-square}, choose the beta code $x$ even.
In true witnesses, $F=t!$, $V=(F^2+1)^t$, and $y=V(V-1)$ is even.
Every beta CRT modulus $1+(i+1)y$ is odd; adding their odd common modulus
changes the code's parity. A sufficiently large even representative
preserves all endpoint and nonnegative quotient requirements.

Each exact factorial-product factor $1+F^2/i$ is odd, since
$F^2/i=F(F/i)$ is even for $1\le i\le t$. Hence $B$, and then $N=B^x$,
are odd. Each of the nine independent CRT representatives can be shifted
by the odd modulus $N^2$, so their parities can be prescribed independently.
Choose opposite parities for $a,b$ as required by the semantic anchor,
and choose one of the other seven parities so that their total sum is even.
Then $H_0$ is odd. All ordinary representatives can retain their old bounds
at most $2N^2$; choose the dominant $a$ in $[2N^2,4N^2)$ with its prescribed
parity. The inherited dominance and frame estimates are therefore preserved.
Finally $BH_0$ is odd and $x$ is even, so
\[
 z=\frac{(BH_0)^{x/2}-1}{2}\in\Z
\]
supplies the odd-square row. This is a construction of fresh witnesses,
not a claim that every old zero already has an odd-square radix.

\subsection{One normalized budget across all four denominator levels}
The incoming Common remainders construction uses a stronger sharing step.
Set $M=q^x$ and let $L_i$ be each remainder's original denominator.
By \eqref{eq:endpointbound}, every exponent defining $L_i$ is at most $x$;
hence $L_i$ is a positive divisor of $M$. For the three smaller denominator
levels introduce one ratio coordinate each, with a quadratic row
\[
 L v_L=M.
\]
For the largest level take $v_M=1$ without adding a coordinate. Replace
an old digit identity $U_i=L_iQ_i+r_i$ by
\begin{equation}\label{eq:normalized-digit}
 v_{L_i}U_i=MQ_i+R_i.
\end{equation}
All displayed high quotients are existing registers or products arranged as
in the literal quadratic graph. In particular, a range digit's old product
$NL_i$ can be replaced by the shared product $NM$, so the new actual row is
$v_{L_i}U_i=(NM)H_i+R_i$. The factorial digit uses the existing register
for $B+qH$. No cubic product is silently inserted into a residual.

\begin{lemma}\label{lem:normalized}
The scaled identities \eqref{eq:normalized-digit}, the ratio rows, and
\begin{equation}\label{eq:one-budget}
 M^2-2\sum_{i=1}^{11}R_i^2=u^2+v^2+w(w+1)
\end{equation}
recover every old remainder interval. Conversely, canonical witnesses with
$0\le r_i<L_i/8$ extend to these rows.
\end{lemma}
\begin{proof}
Each ratio equals $M/L_i>0$. Since both $v_{L_i}U_i$ and $MQ_i$ are
divisible by $v_{L_i}$, so is $R_i$. Thus
$r_i=R_i/v_{L_i}$ is an integer; this is a proof-only division, not a gate
or an extra rational witness. Dividing the scaled row restores the original
digit identity. Equation \eqref{eq:one-budget} gives $|R_i|\le M/\sqrt2<3M/4$, hence
$|r_i|<3L_i/4$. By \eqref{eq:smalltail} and congruence, the difference
$r_i-\rho(a_i,k_i)$ is divisible by $L_i$ and has absolute value below
$7L_i/8<L_i$. Hence it is zero. All deleted interval slacks and product
witnesses can be restored exactly as before.

Conversely, define $v_{L_i}=M/L_i$ and $R_i=v_{L_i}r_i$ from canonical
witnesses. Every scaled row holds, and $0\le R_i<M/8$. Therefore
\[
 2\sum_{i=1}^{11}R_i^2<\frac{11}{32}M^2<M^2.
\]
The integer $M^2-2\sum R_i^2$ is nonnegative and has the required
ternary representation. This proves existence in both directions.
\end{proof}
The incoming common-denominator source instead uses
$M^2-4\sum R_i^2-1=G$. Its soundness radius is $M/2$, and the same
canonical witnesses complete it because $4\sum R_i^2<11M^2/16<M^2$.
The two budgets are therefore existentially equivalent through the calibrated
source relation, not by an asserted pointwise equality of their private
ternary witnesses.

The new form combines the common-denominator idea with the earlier pre-factorial
margin and, where specified, the compatible odd-square radix. In
particular, the ratio rows are interpreted only after positive powers have
been established; the large common guard is never used to establish the
signs on which its own interpretation depends.

\subsection{Certified integral changes preserve the entire zero set}
\begin{lemma}\label{lem:elementary}
Let $R=(R_1,\ldots,R_m)$ be integer polynomials. Each of the following
operations preserves the indicated integer zero set: an integral elementary
coordinate shear; an elementary integral residual-row operation; and
elimination of a coordinate by a monic defining row. The first two are
bijections on the complete zero sets; the third is a bijection after restoring
the eliminated coordinate by its defining polynomial.
\end{lemma}
\begin{proof}
For a shear $x_i=y_i+c y_j$ with $i\ne j$ and $c\in\Z$, leave every
other coordinate unchanged. The inverse is $y_i=x_i-cx_j$, again integral.
Substitution gives $R(x)=R(Ly)$ identically. Translation by an integer
constant has the equally explicit inverse translation.

For a row operation replace $R_i$ by $R_i+cR_j$, with $i\ne j$, keeping
all other rows. A zero of the old rows is a zero of the new rows. A zero of
the new rows has $R_j=0$, hence $R_i=(R_i+cR_j)-cR_j=0$. Equivalently,
the residual vector is multiplied by an elementary matrix in
$\operatorname{GL}_m(\Z)$. Its sum of squares generally changes, but its
zero set does not.

For a row $x_i-f(x_{\ne i})=0$ with $f\in\Z[x_{\ne i}]$, substitute
$x_i=f$ in every remaining row. A new zero extends uniquely by this value,
and every old zero has that value. This proves both directions and
uniqueness. An exact checker verifies that any degree restriction survives
the substitution; it is not a general property of monic elimination.
\end{proof}
A frozen sequence of these operations is a finite mathematical certificate.
The particular heuristic that found it is immaterial to its validity.
A shear transported between two sources by coordinate names requires
checking the full named substitution on the actual target rows; positional
index agreement alone is insufficient. Residual height and support are
measured only after collecting the whole final sum of squares.

\subsection{The exact five selected endpoints}
The starting literal array is the incoming centered Common remainders
252-variable, 182-row equation. Its translation certificate identifies the
semantic radix coordinate. The new construction deletes exactly the old
front guard on $q-B-1$, adds the odd-square row, and changes the common
budget from coefficient four with constant $-1$ to coefficient two with
constant zero. The two witness arguments above justify these semantic
changes. The discarded private guard coordinates are absent afterward;
all other rows are retained exactly at this stage. The resulting source
has 250 variables and 182 rows.

The balanced construction first applies 76 fixed integral coordinate shears,
transported by coordinate identity from the Shared bounds certificate.
Five elementary residual additions expose three affine defining rows.
In the intermediate 250-coordinate order, with zero-based indices and
coordinates $\xi_i$, their unique restorations are
\begin{align*}
 \xi_{37}&=-4-7\xi_{10}-4\xi_{11}-4\xi_{21}+\xi_{22},\\
 \xi_{89}&=-12+6\xi_3+8\xi_{84}+4\xi_{85}+\xi_{86},\\
 \xi_{166}&=-3+2\xi_{107}-2\xi_{148}+\xi_{164}.
\end{align*}
The defining rows have unit coefficients on the eliminated coordinates.
Substituting these formulas and deleting those rows reduces the system to
247 variables and 179 rows. After reindexing, 33 further elementary
residual additions give the intermediate \file{balanced247} checkpoint. Every identity is certified
against the actual preceding stage, and every remaining residual is still
quadratic.

Three further monic quadratic rows can be eliminated without raising
residual degree. In the 247-coordinate checkpoint order, write the coordinates
as $\eta_i$. Their exact restorations are
\begin{align*}
 \eta_{62}&=10+14\eta_{53}+4\eta_{53}^2,\\
 \eta_{64}&=-\eta_{12}+\eta_{12}^2+4\eta_{12}\eta_{63}+4\eta_{63}^2,\\
 \eta_{45}&=2\eta_{10}-\eta_{10}^2-4\eta_{10}\eta_{43}
            +2\eta_{43}-4\eta_{43}^2+\eta_{44}.
\end{align*}
These are the semantic coordinates named \file{mul.281}, \file{add.292}
and \file{mul.194}, respectively. At each deletion the coordinate occurs
only as a linear monomial in every other residual. Such a residual has the
form $av+G$ with $a\in\Z$ and $G$ quadratic and independent of $v$;
replacing $v$ by its quadratic restoration keeps degree at most two.
The independent checker verifies this occurrence condition before every
quadratic deletion, in addition to checking all resulting degrees. Unique
restoration proves the same zero-set bijection as Lemma~\ref{lem:elementary}.
After these three deletions and nine further row additions, the result has
244 variables, 176 rows, 5212 terms and height 3968.

The 243-variable alternative also eliminates the coordinate \file{add.380}
by a 15-term affine restoration, before those three quadratic deletions.
In the same checkpoint order that additional formula is
\begin{align*}
 \eta_{103}={}&5\eta_{10}-3\eta_{11}+4\eta_{21}+2\eta_{24}
 +4\eta_{26}+2\eta_{29}+2\eta_{31}+2\eta_{34}\\
 &+2\eta_{36}+4\eta_{38}+2\eta_{39}+2\eta_{42}
 +2\eta_{46}+4\eta_{67}+2\eta_{68}.
\end{align*}
Its final frozen residual basis uses 802 further row additions. The sparse
form instead uses the original 76 shears and seven residual additions,
without any projection. The separate 246-variable form uses six elementary
linear-coordinate shears, 1057 row additions and four affine projections.
The low-height form retains the first three affine projections and a longer
frozen row basis. No optimization search is needed to define or reproduce
any of these finite sequences.

\begin{center}\small
\begin{tabular}{@{}lrrrr@{}}\toprule
Endpoint & Variables & Rows & Terms & Height\\\midrule
\file{balanced244} &244&176&5212&3968\\
\file{dimension243} &243&175&19517&7680\\
\file{sparse250} &250&182&4770&6709\\
\file{dimension246} &246&178&21537&6144\\
\file{lowheight247} &247&179&10422&3840\\\bottomrule
\end{tabular}
\end{center}
Every equation is the sum of the squares of its full residual array over
$\Z$. The first four contain the collected monomial $13x_0^4$; the last
contains $22x_0^4$. Since every row has degree at most two, these nonzero
coefficients establish exact degree four. Full coefficient collection gives
the displayed support and height, rather than an upper bound from the
largest residual coefficient. The intermediate \file{exposed247} and
\file{balanced247} arrays remain technical source and construction
checkpoints; they are not presented as additional selected frontier points.

The inverse algebraic witness maps are explicit: reverse the residual
operations, restore each eliminated coordinate, and reverse the coordinate
shears. For the two semantic source changes, restore old guards with fresh
ternary witnesses after the proven residue and sign recovery. Conversely,
start with a finite refutation, use the compatible even-code and odd-radix
canonical witnesses, fill the common spherical budget, then apply the
exact coordinate maps. This composes both existential directions all the
way to the retained calibrated source. The semantic stages are not claimed
to give a bijection between arbitrary private witness choices.

